\subsection{Simple Algebras over Finite Fields}

THEOREM 2 (Wedderburn). --- Every finite field is commutative.

Let D be a finite field, and let K be its center. The K-algebra D is central simple of degree \(m^{2}\) , where m is a strictly positive integer. The reduced norm is a homogeneous polynomial mapping \(Nrd : D \to K\) of degree m (VIII, p.~345, Proposition 6), and we have \(\text{Nrd}(a) \neq 0\) for every \(a \neq 0\) in D (VIII, p.~340, Proposition 3). The corollary above implies that \(m \geqslant m^{2}\) , and therefore m = 1. So we have D = K.

COROLLARY 1. --- Every finite simple ring is isomorphic to a matrix ring \(\mathbf{M}_{n}(\mathrm{L})\) , where n is a strictly positive integer and L is a finite commutative field.

This follows from Theorem 2 and the structure theorem for simple rings (VIII, p.~120, Theorem 1).

COROLLARY 2. --- Let K be a finite commutative field. Every central simple algebra over K is isomorphic to a matrix algebra \(\mathbf{M}_{n}(\mathrm{K})\) , where n is a strictly positive integer.

This follows from Theorem 2 and the structure theorem for central simple algebras (VIII, p.~252, Theorem 1).

Remarks. --- 1) Here is another proof of Theorem 2. Let D be a finite field, let K be its center, and let L be a maximal commutative subfield of D. Let x be an element of \(D^{*}\) . It belongs to a maximal commutative subfield \(L_{1}\) of D. We have the equality

\[
[ \mathrm{D}: \mathrm{K} ] = [ \mathrm{L}: \mathrm{K} ] ^ {2} = [ \mathrm{L} _ {1}: \mathrm{K} ] ^ {2}
\]

by Corollary 2 of VIII, p.~265, and therefore \([\mathrm{L}:\mathrm{K}] = [\mathrm{L}_1:\mathrm{K}]\) . By Proposition 3 of V, §12, No.~2, p.~94, the extensions L and \(\mathrm{L}_1\) of K are isomorphic. By VIII, p.~263, Corollary, there exists an element \(a\) of \(\mathrm{D}^*\) such that \(a\mathrm{La}^{-1} = \mathrm{L}_1\) , so \(a^{-1}xa\) belongs to L. We then have \((ay)^{-1}x(ay) = a^{-1}xa\) for every \(y\in \mathrm{L}^*\) . Consequently, if S is a set of representatives of the left cosets of \(\mathrm{D}^*\) modulo \(\mathrm{L}^*\) , then every element of \(\mathrm{D}^* -\{1\}\) can be written as \(sxs^{-1}\) with \(s\in \mathrm{S}\) and \(x\in \mathrm{L}^* -\{1\}\) . We denote the order of \(\mathrm{D}^*\) by \(d\) and that of \(\mathrm{L}^*\) by \(\ell\) . Since the cardinal of S is equal to \(d / \ell\) , we have \(d - 1\leqslant (d / \ell)(\ell -1) = d - d / \ell\) . It follows that \(\ell = d\) , and therefore \(\mathrm{L} = \mathrm{D}\) , which proves that the field D is commutative.

\begin{enumerate}
\def\labelenumi{\arabic{enumi})}
\setcounter{enumi}{1}
\tightlist
\item
  Let L be a commutative field with the following property:
\end{enumerate}

(C \(_{1}\) ) Let V be a vector space of finite dimension n over the field L, and let d be such that 0 \textless{} d \textless{} n.~For every homogeneous polynomial mapping F : V → L of degree d, there exists a nonzero element x of V such that F(x) = 0.

The proof of Theorem 2 shows that every field of finite degree over L with center L is equal to L.

By the corollary of VIII, p.~356, every finite field has property \((\mathrm{C}_{1})\) . We can prove (VIII, p.~360, Exercise 7) that the following fields have property \((\mathrm{C}_{1})\) :

- every algebraic extension of a finite field

-- every field of rational fractions in one variable with coefficients in an algebraically closed field (Tsen's theorem).

- *every field endowed with a discrete valuation for which it is complete and whose residue field is algebraically closed (VIII, p.~332, Exercise 17).*

\begin{enumerate}
\def\labelenumi{\arabic{enumi})}
\setcounter{enumi}{2}
\tightlist
\item
  Suppose that the field \(\mathbf{K}\) satisfies the following condition:
\end{enumerate}

- If \(\mathrm{L}\) is an extension of \(\mathrm{K}\) of finite degree, then it is cyclic and the norm mapping \(\mathrm{N}:\mathrm{L}^{*}\to \mathrm{K}^{*}\) is surjective.

This condition is satisfied, in particular, when the field K is finite (V, §12, No.~2, p.~95, Proposition 4). We can then prove that every field of finite degree over K with center K is equal to K (Exercise 10 of VIII, p.~329).

